\noindent\hspace*{2em}神经网络计算图的多核执行能够提高计算并行度，而任务切分带来的跨域搬运、片上容量占用与共享带宽竞争会影响总体工期。

本文研究独立任务、同核复用与共享只读缓存三种场景下的多核切图与调度。首先通过\textbf{计算依赖提取与弱连通分量分析}，构造子图划分与核心分配方案；接着根据张量的消费范围和物理副本生存区间，建立结构搬运与容量约束模型；最后引入 \textbf{FIFO 缓存事件模型}，由题给 Python 事件模拟程序评价 100 张计算图在 1 至 5 核下的调度结果。

\textbf{针对问题一}，建立 \textbf{Task 驻留域结构搬运模型}。首先提取计算依赖及张量的生产、消费关系，按接收 Task 统计输入读入、跨任务搬运与最终写回，并根据同核次序和前驱依赖确定 Task 激活条件。接着以\textbf{整图分配、分量装箱与拓扑切块}生成候选，用\textbf{静态评分}安排评价顺序，再由事件评估程序按 Makespan 优先、额外搬运次之选定方案。相对整图单核基准，五核平均速度比为 \textbf{3.4383}；五核初始方案到最终方案的平均工期改善为 \textbf{58.6697\%}，\textbf{89 张图}得到改善。

\textbf{针对问题二}，建立\textbf{核心驻留与物理副本容量模型}。首先将同核子图合并为核级 Task，按接收核心统计跨核搬运。其次根据物理副本从空间申请至最后读者完成的存活区间，检查 L1、UB 的同时占用，分别核算结构搬运与换出恢复。然后以\textbf{关键路径优先}和释放优先生成拓扑切块候选，结合贪心分核与局部调整，经题给事件评价选定方案。五核平均速度比为 \textbf{3.4738}；初始方案到最终方案的平均工期改善为 \textbf{60.3622\%}，\textbf{98 张图}得到改善。

\textbf{针对问题三}，建立\textbf{基于查询与完成事件的 FIFO Cache 模型}。首先按照发射时查询、读取完成后插入和先进先出淘汰的规则，把命中读与未命中读分别纳入 Cache 和 DDR 的独立带宽服务。接着以问题二最终计划为起点，通过共享输入聚合、迁块和核内调序构造候选，并以\textbf{同图同核配对}区分硬件作用、调度适配和综合效果。五核平均速度比为 \textbf{3.5533}；三段效应比分别为 \textbf{1.0146}、\textbf{1.0139} 和 \textbf{1.0289}。

最后，进行\textbf{受控消融与灵敏度分析}。六张代表图的 Cache 容量与带宽分别减半、加倍，\textbf{24 组}固定计划全部通过官方评价；在硬件变体下重新生成的 \textbf{37 条}候选也全部通过，每种变体均有 \textbf{3/6} 组改变计划。重优化相对默认配置的逐图工期比均值分别为容量减半 \textbf{0.9900}、容量加倍 \textbf{0.9869}、带宽减半 \textbf{0.9892}、带宽加倍 \textbf{0.9885}。

\noindent{\zihao{3}\textbf{关键词：}}多核调度；张量驻留域；拓扑切块；容量约束；FIFO 缓存
